\subsection{The retained multiplier and every even internal root}\label{proof:R2-C042}
\paragraph{R42.1 Original model and precise claims}

For a real skew kernel K with pointwise absolute value at most one on
one original finite probability space, put
\fitdisplay{
B=-K,\quad u=B\one,\quad T=J+B,\quad A=J-B=T^*,\quad
 p=1-u,\quad q=1+u,\quad d=E u^2\in[0,1/3].
}{}
All adjoints, norms and expectations use this unchanged original measure.
Write \(f_n=A^{n}\one,\allowbreak  \bar f_n=T^{n}\one,\allowbreak  h_n=T f_{n-1},\allowbreak  W_n=A \bar f_{n-1}. M_v\) denotes multiplication by the actual function v. After zero-mass
atoms are removed, p(x), \(q(x)\ge \mu(x)>0\); there is no uniform cutoff.
Define \(I_q(V)=\E[q^{-1/3}|V|^{8/3}]\).

Theorem A — arbitrary real fourth input with an actual arm.
For every real original Y, every \(a\ge 2\) and every \(M\ge 1,\allowbreak \)
\fitdisplay{
I_q(T^M(\bar f_aY))^3
 \le(1-d/2)(1-d/72)^2\,N_{M-1}^4(E|Y|^4)^2,
}{R42.1.1}
where valid original norm bounds are
\fitdisplay{
N_0=1,\quad N_1=1+32d/19,\quad
 N_2=(1+(1024/855)\delta_2)^2,
}{}
\fitdisplay{
N_n=(1-d/8)^{\lfloor n/3\rfloor}\quad(n\ge3),
 \qquad \delta_2=\|B^21\|^2\le(13/32)d.
}{R42.1.2}
Consequently a useful all-length exponential form is
\fitdisplay{
I_q(T^M(\bar f_aY))^3\le e^{c_Md}(E|Y|^4)^2,
\quad
 c_M=\begin{cases}
 -19/36,&M=1,\\
 4247/684,&M=2,\\
 11507/3420,&M=3,\\
 -19/36-\lfloor(M-1)/3\rfloor/2,&M\ge4.
 \end{cases}
}{R42.1.3}
In particular \(c_M\le 4247/684\) for EVERY M. Thus any actual H with
\(EH^4\le \exp(-\lambda\) d) has the same-direction output
\fitdisplay{
I_q(T^M(\bar f_aH))^3
 \le\exp[-(2\lambda-c_M)d]
 \le\exp[-(2\lambda-4247/684)d].
}{R42.1.4}
No ordinary eighth hypothesis occurs in this theorem.

Theorem B — every even root with at least four residual arms.
At an original internal root let the maximally unpaired side-arm lengths
be strictly increasing \(r_{1}<...<r_{2j},\allowbreak  j\ge 2.\) Its terminal has any finite
strict residual list of size at least two. All original equal-pair
decorations, all independent \(innerL\ge 1\) and \(outerM\ge 1\) are allowed.
Section3 specifies all original directions from these integers alone.
The actual output \(U_D\) satisfies
\fitdisplay{
I_q(U_D)^3\le
 \begin{cases}
 e^{-15701d/3420},&r_1=1,\\
 e^{-3277d/684},&r_1\ge2.
 \end{cases}
}{R42.1.5}
The common first coefficient applies to the whole family.

Theorem C — complete unit-minimum even roots and original P5.
Every unit-minimum even residual, including an exact pair(1, t), has,
for every nonreducible terminal, an all-L/M actual output satisfying
\(Iq(U)^3\le \exp(-d/72).\)
For a P5 whose two noncentral internal roots each either have an even
residual of size at least four, or have an exact pair(1, t), the two
terminal and central multisets may be arbitrary, and all four positive
core lengths independent: there is a structure-only direction with \(Z\le 1.\)
In the nonreducible-end, central-R branch, the common deficit is d/288.
If both internal residuals have size at least four, it improves to
15701d/13680. Section5 gives precisely qualified mixed-side consequences.

\paragraph{R42.2 The complete analytic proof: retain and pay the multiplier}

\paragraph{R42.2.1 Exact original ingredients}

The actual ordinary multiplier statement is
\fitdisplay{
\|T M_{\bar f_a}\|_{2\to2}^2\le1-d/72\quad(a\ge2).
}{R42.2.1}
This is C064(1.6)/C066(1.4). Q4/C115(8.3) states
\(\lVert M_{f_a}T\rVert ^2\le 1-d/72.\) Replacing K by -K changes that operator to
\(M_{\bar f_a}A\); taking its ordinary adjoint gives T \(M_{\bar f_a},\allowbreak \) exactly (R42.2.1).
Reversal and adjunction use the same probability measure.
This estimate uses the full multiplier.

C100 TheoremA (A068) gives
\fitdisplay{
E[\bar f_n^4/q]\le1-(13639/23040)d\le1-d/2\quad(n\ge4).
}{R42.2.2}
The distinct n3 bound is the C056 inequality
E[\(\bar f_3^4/q\)]\ensuremath{\le}\(1-23d/36\le 1-d/2.\) Thus (R42.2.2)'s weaker final bound holds for
EVERY \(n\ge 3.\) The C100 proof in Sections2-3 uses
two retained finite convolution losses, a truncated quartic majorant and
the true correlated cubic. The derivation uses these finite bounds directly.

The ordinary T power bounds are exactly (R42.1.2). N2 is the separate
two-step bound; the \(n\ge 3\) bound uses \(separate3/4/5\) bases and finite
blocks. In particular no assumption \(\lVert T^{n+1}\rVert \le \lVert T^n\rVert \) is introduced.

\paragraph{R42.2.2 Every real input, every actual positive path length}

Let P be any original positive-kernel operator and \(h=P1.\) Pointwise row
Cauchy and original Holder with exponents3 \(and3/2\) give
\fitdisplay{
|PY|^2\le hP(|Y|^2),
}{}
\fitdisplay{
I_q(PY)\le
 \{E[h^4/q]\}^{1/3}
 \{E[(P(|Y|^2))^2]\}^{2/3}.
}{}
The ordinary L2 norm then yields
\fitdisplay{
I_q(PY)^3\le E[h^4/q]\cdot\|P\|^4(E|Y|^4)^2.
}{R42.2.3}
This calculation uses \(|Y|^2,\allowbreak \) so it allows any real Y, without a sign or
actual-tree assumption on Y. The original positive row h is unchanged.

Now put \(P=T^M M_{\bar f_a}.\) Its EXACT row and factorization are
\fitdisplay{
h=T^M \bar f_a=\bar f_{M+a},\qquad
 P=T^{M-1}(T M_{\bar f_a}).
}{R42.2.4}
Because \(M+a\ge 3,\allowbreak \) the bare-row bound applies at every allowed boundary,
including a2, M1. The actual norm satisfies
\fitdisplay{
\|P\|^4\le\|T^{M-1}\|^4(1-d/72)^2.
}{R42.2.5}
Inserting (R42.2.2),(R42.2.5) into (R42.2.3) proves (R42.1.1). In particular we did NOT
replace \(\lVert T^M M_{\bar f_a}\rVert \) by \(\lVert T^M\rVert .\) The actual multiplier is paid before
the remaining M-1 original edges; all operator orders are retained.

For (R42.1.3), \(\log(1-d/2)+2\log(1-d/72)\le -(19/36)d.\)
When M2 the remaining norm contributes at most128d/19, and
\(128/19-19/36=4247/684.\) When M3 it contributes at most
\(8(1024/855)\delta_2\le 3328d/855,\allowbreak \) and
\(3328/855-19/36=11507/3420.\) For \(M\ge 4,\allowbreak \) (R42.1.2) contributes at most
-floor((M-1)/3)d/2. M1 has no remaining norm cost. These four exact
cases cover every positive integer M and prove the common maximum.

At d0, the actual \(\bar f_a\) and q equal1 although B need not vanish.
The same norm and positive-operator proof gives (R42.1.1) for all realY.
Leaf integration is used only later for actual tree messages, not for an
arbitrary real function. No division by d occurs in the transfer proof.

\paragraph{R42.3 Actual orientation algorithm for every even residual of \(size\ge 4\)}

Fix an original ordering of equal-length arms to choose their maximal
pairing. Orient each real equal pair as a valid \(f_l h_l\) pair or its
complete reversal. Assign ALL root equal pairs to the inner factor H;
terminal pairs remain in \(E_t\) at the original terminal. Positivity permits comparison while leaving the original edges and $e(U)$ unchanged.

\paragraph{R42.3.1 Minimum one: omit only the second arm from H}

Write the actual root list as
\fitdisplay{
(1,t,a,b,s_1,v_1,\ldots,s_k,v_k),\quad
 1<t<a<b<s_1<v_1<\cdots,
 \qquad k=j-2\ge0.
}{R42.3.1}
Retain the actual t arm fully inward, as \(\bar f_t.\) Every remaining root arm
forms the actual odd list(1, a, b, s1, v1,...), with \(a\ge 3.\) Use EXACTLY C186
Section1 at that root and terminal. The root factor is
\fitdisplay{
H=E_r pW_a \bar f_b\prod_i W_{s_i}\bar f_{v_i}
          T^L(E_tF_{\Sigma,L}^{186}),\qquad
 U=T^M(\bar f_tH).
}{R42.3.2}
The C186 ordinary fourth gives \(EH4\le \exp(-27d/5)\) for allk, L.
Its later-pair improvement is optional; the base suffices. Apply (R42.1.4):
\fitdisplay{
2(27/5)-4247/684=15701/3420>0.
}{R42.3.3}
This includes t2, a3, b4 and every finite post list. No minimum-four post
restriction is imposed on the first actual pair \(W_a \bar f_b:\) C186 itself
includes a3. The \(\bar f_t\) arm is not inside H and is used exactly once outside.

\paragraph{R42.3.2 Minimum at least two: omit only the shortest arm from H}

Write the original root list as
\fitdisplay{
(s,a,b,c,s_1,v_1,\ldots,s_k,v_k),\quad
 2\le s<a<b<c<s_1<v_1<\cdots,\quad k=j-2\ge0.
}{R42.3.4}
Retain the actual s arm fully inward, as \(\bar f_s.\) The remaining actual
odd root(a, b, c, s1, v1,...) has \(a\ge 3\) and falls exactly in C189 TheoremA.
Define \(Pplus=h_b f_c\) product \(h_si f_{v_i}\) and
\(Pminus=W_b \bar f_c\) product \(W_si \bar f_{v_i}.\) The actual root/terminal choices are:


\begingroup\small
\setlength{\tabcolsep}{3pt}
\setlength{\LTleft}{0pt}\setlength{\LTright}{\fill}
\begin{longtable}{@{}p{0.2200\linewidth}p{0.3500\linewidth}p{0.3500\linewidth}@{}}
\caption{The retained multiplier and every even internal root: parameter values.}\label{tab:parameters:r2-c042:1}\\
\toprule
Entry & Actual root factor inside H & Actual F \\
\midrule\endfirsthead
\toprule
Entry & Actual root factor inside H & Actual F \\
\midrule\endhead
\bottomrule\endfoot
Exact(1, z), L1 & \(f_a Pminus\); & \(p\bar f_z\). \\
\addlinespace[2pt]
Exact(1, z), \(L\ge 2\) & \(W_a Pminus\); & \(q f_z\). \\
\addlinespace[2pt]
Every other terminal & \(f_a Pplus\); & The exact C189 table below. \\
\end{longtable}
\endgroup


In all cases \(H=E_r(rootfactor)T^L(E_tF),\allowbreak  U=T^M(\bar f_sH).\)
C189's ordinary fourth is \(EH4\le \exp(-11d/2),\allowbreak \) including both
special(1, z) rows and all equal pairs. The new output coefficient is
\fitdisplay{
11-4247/684=3277/684>15701/3420.
}{R42.3.5}
The special first root arm \(W_a\) is not changed to \(f_a\) by notation.
The negative pairs in the two short-one rows are not changed to Pplus.
The s arm is distinct from all displayed H arms and appears only as \(\bar f_s.\)

\paragraph{R42.3.3 The complete inherited terminal directions used in both rows}

The following table specifies the terminal directions. All
lists are strictly increasing and later products contain only displayed
original pairs. The two ordinary-root bounds agree here except for the
exact(1, z), L1 row described above.


\begingroup\small
\setlength{\tabcolsep}{3pt}
\setlength{\LTleft}{0pt}\setlength{\LTright}{\fill}
\begin{longtable}{@{}p{0.2200\linewidth}p{0.3500\linewidth}p{0.3500\linewidth}@{}}
\caption{The retained multiplier and every even internal root: parameter values.}\label{tab:parameters:r2-c042:2}\\
\toprule
Entry & C186 F & C189 F \\
\midrule\endfirsthead
\toprule
Entry & C186 F & C189 F \\
\midrule\endhead
\bottomrule\endfoot
Exact(1, z), \(z\ge 2,\allowbreak \) L1 & \(q f_z\); & \(p\bar f_z\). \\
\addlinespace[2pt]
Exact(1, z), \(z\ge 2,\allowbreak  L\ge 2\) & \(q f_z\); & \(q f_z\). \\
\addlinespace[2pt]
Exact(2, 3) & h2 f3; & h2 f3. \\
\addlinespace[2pt]
Exact(2, z), \(z\ge 4\) & \(\bar f_2 W_z\); & \(\bar f_2 W_z\). \\
\addlinespace[2pt]
Exact(3, 4) & h3 f4; & h3 f4. \\
\addlinespace[2pt]
Exact(3, z), \(z\ge 5,\allowbreak \) L1 & \(h_{3} f_z\); & \(h_{3} f_z\). \\
\addlinespace[2pt]
Exact(3, z), \(z\ge 5,\allowbreak  L\ge 2\) & \(\bar f_3 W_z\); & \(\bar f_3 W_z\). \\
\addlinespace[2pt]
Single pair \(minimum\ge 4\) & \(h_\alpha f_\beta\); & same. \\
\addlinespace[2pt]
Every even list with at least two pairs & \(\prod h_{\alpha_i} f_{\beta_i}\); & same. \\
\addlinespace[2pt]
Exact(1, A, C), \(2\le A<C\) & \(q f_A h_C\); & same. \\
\addlinespace[2pt]
Unit-minimum odd multipair with(A1, C1, \(A2)=(2,\allowbreak \) 3, 4) & \(q f_{2} h_{C_2} h_{3} f_{4} \prod_{i\ge 3}h_{A_i} f_{C_i}\); & same. \\
\addlinespace[2pt]
All other unit-minimum odd multipairs & \(q f_{A_1} h_{C_1} \prod_{i\ge 2}h_{A_i} f_{C_i}\); & same. \\
\addlinespace[2pt]
Nonunit-minimum odd list, first paired \(shortB1=3\) & \(\bar f_2 f_{3} h_{C_1} \prod_{i\ge 2}h_{B_i} f_{C_i}\); & same. \\
\addlinespace[2pt]
Nonunit-minimum odd list, first paired \(shortB1\ge 4\) & \(\bar f_{A} \prod_i h_{B_i} f_{C_i}\); & same. \\
\end{longtable}
\endgroup


The first paired \(shortB1=3\) forces \(A=2.\) The crossed unit row requires at
least five residual arms. Exactly three unit-minimum arms use the triple row. This gives the stated estimate for every finite nonreducible terminal. All inner connector
edges are T and all outer connector edges T; the arm turns are exactly
the h/W definitions. 

\paragraph{R42.4 Complete unit-minimum even roots, including the pair boundary}

An even residual starting at one either has at least four entries and
uses (R42.3.1)-(R42.3.3), or is exactly(1, t), \(t\ge 2.\) In the latter case use the
complete Q5-C106 ordinary table \(F_ord\) for its actual nonreducible terminal:
\fitdisplay{
H=E_r pT^L(E_tF_{ord}),\quad EH8\le1-d/20000,
 \quad U=T^M(\bar f_tH).
}{}
The original C066 arbitrary-real eighth theorem yields
\fitdisplay{
I_q(U)^3\le(1-d/72)(1-d/20000)\le e^{-d/72}.
}{R42.4.1}
The exact-pair boundary uses the eighth-moment theorem for the C106 rows.
The structural branches have their stated directions. Together they prove
the first assertion of TheoremC. TheoremB's nonunit minimum rows also
have a much stronger output than exp(-d/72).
